\documentclass[a4paper]{article}
% Español (México) con XeLaTeX: fontspec para fuentes Unicode, polyglossia
% para la división silábica y los nombres del documento en la variante mexicana.
\usepackage{fontspec}
\usepackage{polyglossia}
\setdefaultlanguage[variant=mexican]{spanish}
% Los nombres chinos van como «pinyin (original)»: xeCJK da a los caracteres
% Han su propia fuente; sin ella XeLaTeX los omite con un aviso.
% FandolSong no cubre algunos caracteres de nombres (U+73FA); WenQuanYi Zen Hei sí.
\usepackage[AutoFallBack=true]{xeCJK}
\setCJKmainfont{FandolSong-Regular.otf}
\setCJKfallbackfamilyfont{\CJKrmdefault}{WenQuanYi Zen Hei}
% polyglossia no traduce los nombres de listings.
\AtBeginDocument{\providecommand{\lstlistingname}{}\renewcommand{\lstlistingname}{Listado}%
  \providecommand{\lstlistlistingname}{}\renewcommand{\lstlistlistingname}{Índice de listados}}
\usepackage{amsmath,amssymb}
\usepackage{graphicx}
\usepackage[margin=2.35cm]{geometry}
\usepackage[most]{tcolorbox}
\usepackage{etoolbox}
\usepackage{listings}
\usepackage{xcolor}
\usepackage{hyperref}
\usepackage{booktabs,longtable,tabularx,array}
\usepackage{subcaption}
\usepackage{float}
\usepackage{enumitem}
\usepackage{microtype}
\usepackage{fancyhdr}
\usepackage{needspace}

\definecolor{kbBack}{HTML}{F2F6FA}
\definecolor{kbFrame}{HTML}{9DB4CC}
\definecolor{kbTitle}{HTML}{52708F}
\definecolor{ibBack}{HTML}{FBF5E7}
\definecolor{ibFrame}{HTML}{C9A961}
\definecolor{ibTitle}{HTML}{7D5F1E}
\definecolor{wbBack}{HTML}{FCF2F0}
\definecolor{wbFrame}{HTML}{D6A096}
\definecolor{wbTitle}{HTML}{9E4F43}
\definecolor{dbBack}{HTML}{F4F7F3}
\definecolor{dbFrame}{HTML}{AEC2AC}
\definecolor{dbTitle}{HTML}{5F7F64}

\tcbset{softbox/.style={
  enhanced,breakable,boxrule=0.8pt,arc=2pt,left=3mm,right=3mm,
  top=2mm,bottom=2mm,coltitle=white,fonttitle=\bfseries,
  toptitle=0.6mm,bottomtitle=0.6mm
}}
\newtcolorbox{knowledgebox}[1]{softbox,colback=kbBack,colframe=kbFrame,colbacktitle=kbTitle,title={#1}}
\newtcolorbox{importantbox}[1]{softbox,colback=ibBack,colframe=ibFrame,colbacktitle=ibTitle,title={#1}}
\newtcolorbox{warningbox}[1]{softbox,colback=wbBack,colframe=wbFrame,colbacktitle=wbTitle,title={#1}}
\newtcolorbox{dialoguebox}[1]{softbox,colback=dbBack,colframe=dbFrame,colbacktitle=dbTitle,title={#1}}

\lstset{
  language=Python,basicstyle=\ttfamily\small,
  keywordstyle=\color{kbTitle}\bfseries,stringstyle=\color{wbTitle},
  commentstyle=\color{dbTitle},breaklines=true,frame=single,numbers=left,
  numberstyle=\tiny\color{gray},captionpos=b,columns=fullflexible,
  keepspaces=true,showstringspaces=false,extendedchars=true
}
\BeforeBeginEnvironment{lstlisting}{\Needspace{12\baselineskip}}

\hypersetup{colorlinks=true,linkcolor=kbTitle,urlcolor=kbTitle,citecolor=wbTitle}
\setlist{itemsep=0.25em,topsep=0.4em}
\setlength{\parindent}{2em}
\setlength{\parskip}{0.25em}
\newcolumntype{L}[1]{>{\raggedright\arraybackslash}p{#1}}

\providecommand{\notetitle}{Notas del curso: [tema del curso]}
\providecommand{\noteauthors}{Elaboradas a partir de los materiales públicos de Stanford CS25}
\providecommand{\notedate}{Versión reescrita en 2026}
\providecommand{\videochannel}{Stanford Online}
\providecommand{\videopublishdate}{2022}
\providecommand{\videoduration}{[duración del video]}
\providecommand{\videourl}{}
\providecommand{\videocoverpath}{cover.jpg}
\providecommand{\coursesite}{https://web.stanford.edu/class/cs25/}
\providecommand{\repourl}{https://github.com/hqhq1025/ai-course-notes}
\providecommand{\skillversion}{wdkns/wdkns-skills@39f1a04}

\newcommand{\figsource}[1]{\par\vspace{-0.3em}{\footnotesize\color{gray}\emph{Fuente: #1}}\par}
\newcommand{\teachervoice}[1]{\begin{knowledgebox}{Nota de clase}#1\end{knowledgebox}}
\newcommand{\term}[1]{\textbf{#1}}

\pagestyle{fancy}
\fancyhf{}
\fancyfoot[C]{\thepage}
\fancyfoot[R]{\footnotesize\color{gray}\href{\repourl}{AI Course Notes}}
\renewcommand{\headrulewidth}{0pt}
\renewcommand{\footrulewidth}{0pt}

\newcommand{\makecscover}{
\begin{titlepage}
\centering
\vspace*{0.7cm}
{\Large Stanford CS25 · Transformers United\par}
\vspace{0.8cm}
{\huge\bfseries \notetitle\par}
\vspace{0.6cm}
{\large \noteauthors\par}
\vspace{0.25cm}
{\large \notedate\par}
\vspace{0.9cm}
\ifdefempty{\videocoverpath}{}{
  \includegraphics[width=0.86\textwidth,height=0.43\textheight,keepaspectratio]{\videocoverpath}\par
}
\vfill
\begin{tcolorbox}[width=0.92\textwidth,colback=black!2!white,colframe=black!45,boxrule=0.7pt,arc=2pt]
\textbf{Autor o canal del video}: \videochannel\par
\textbf{Fecha de publicación}: \videopublishdate\par
\textbf{Duración del video}: \videoduration\par
\textbf{Sitio del curso}: \href{\coursesite}{\nolinkurl{\coursesite}}\par
\ifdefempty{\videourl}{}{\textbf{Enlace del video}: \href{\videourl}{\nolinkurl{\videourl}}\par}
\textbf{Normas de elaboración}: \skillversion\ + las normas source-first / teacher-voice / visual-QA de este repositorio
\end{tcolorbox}
\end{titlepage}
\tableofcontents
\newpage
}


\renewcommand{\notetitle}{Lecture 20: Generalist Agents in Open-Ended Worlds}
\renewcommand{\noteauthors}{Elaborado a partir de la conferencia de Jim Fan, los subtítulos oficiales hechos por personas, las diapositivas recuperadas y artículos como MineDojo / Voyager / Eureka / VIMA}
\renewcommand{\notedate}{CS25 V3 · versión reescrita source-first de 2026}
\renewcommand{\videochannel}{Stanford Online}
\renewcommand{\videopublishdate}{Clase: 2023-10-24; subido: 2023-12-09}
\renewcommand{\videoduration}{47:48}
\renewcommand{\videourl}{https://www.youtube.com/watch?v=wwQ1LQA3RCU}
\renewcommand{\videocoverpath}{cover.jpg}

\newcommand{\lecturefigure}[3]{%
\begin{figure}[H]
  \centering
  \includegraphics[width=0.94\textwidth,height=0.54\textheight,keepaspectratio]{slides-images/#1}
  \caption{#2}
  \figsource{#3}
\end{figure}
}

\let\standardtableofcontents\tableofcontents
\renewcommand{\tableofcontents}{%
  \begingroup
  \footnotesize
  \renewcommand{\baselinestretch}{0.92}\selectfont
  \setlength{\parskip}{0pt}
  \standardtableofcontents
  \endgroup
}

\begin{document}
\makecscover

\section{Auditoría de fuentes: del «gatito pasivo» a los agentes de mundo abierto}

La pregunta central de esta clase no es «cómo construir otro bot de Minecraft más potente», sino cómo conectar los amplios conocimientos previos de un foundation model a un ciclo cerrado capaz de \textbf{observar, actuar, asumir consecuencias y acumular habilidades de forma continua}. Jim Fan primero usa un experimento clásico sobre desarrollo visual para mostrar la importancia de la experiencia activa, y después enlaza MineDojo, MineCLIP, Voyager, Eureka y VIMA en una misma línea de sistemas: el entorno abierto ofrece tareas ilimitadas, internet aporta el conocimiento humano, el foundation model proporciona representaciones transferibles, y reward, code, memory, curriculum y simulator feedback hacen que esos conocimientos previos se traduzcan efectivamente en acción.

\subsection{Límites oficiales de la clase y niveles de evidencia}

Esta reescritura toma como fuente del orden de la clase la grabación oficial de Stanford Online \texttt{wwQ1LQA3RCU}. La clase se impartió el 24 de octubre de 2023, la grabación dura 47 minutos con 48 segundos e incluye subtítulos \texttt{en-US} hechos por personas. Se recuperaron 61 estados didácticos de la grabación en 1080p, que cubren todas las diapositivas que cumplen una función real de concepto, arquitectura, resultado o ejemplo; los fotogramas reproducidos de nuevo y los fotogramas de pura transición no se cuentan dos veces. Los artículos se usan solo para verificar terminología y mecanismos; no sirven para presentar como hechos conocidos en la clase resultados posteriores a octubre de 2023.

\begin{longtable}{@{}L{0.18\textwidth}L{0.31\textwidth}L{0.40\textwidth}@{}}
\toprule
Nivel de evidencia & Material local & Función en las notas \\
\midrule
Grabación oficial & 47:48, 61 estados didácticos & Determina el orden de exposición, la evidencia visual, las demos y los límites de las conclusiones de la clase \\
Subtítulos oficiales hechos por personas & 985 caption, timed transcript & Recuperan motivaciones, advertencias, metáforas, condiciones restrictivas y transiciones entre secciones \\
Primary sources & MineDojo, Voyager, Eureka, VIMA, VPT, RT-2, RoboCat & Verifican objetivos de entrenamiento, nombres de arquitecturas, objetos experimentales y el significado de las fórmulas \\
Auditoría de cobertura y visual & manifest, selection TSV, coverage matrix, PDF QA & Demuestran que se usó cada figura didáctica y que el producto final es legible; no se sustituye la aceptación por «el archivo existe» \\
\bottomrule
\end{longtable}

\begin{warningbox}{Límites temporales y límites de la evidencia}
Este texto trata la generalist-agent research agenda de la clase de 2023. Que MineDojo, Voyager o Eureka obtengan ventaja en cierto simulator y cierta métrica no equivale a haber logrado una inteligencia general que funcione entre entornos y entre embodiments; que VIMA, RT-2 y RoboCat muestren transferencia multimodal y robótica tampoco equivale a haber resuelto la seguridad, la memoria a largo plazo, la evaluación en mundo abierto ni el cuello de botella de los datos del mundo real.
\end{warningbox}

\subsection{El experimento de Held y Hein: por qué la misma entrada visual produce capacidades distintas}

El curso arranca con el experimento de los «dos gatitos» de Held y Hein para elevar la embodiment de «tener un cuerpo robótico» a una condición de aprendizaje más estricta: \term{embodiment (corporeidad)} significa que la percepción y la acción del agente forman un ciclo cerrado a través del cuerpo y de la dinámica del entorno; la acción cambia lo que se verá en el instante siguiente y además produce éxitos, fracasos y costos atribuibles. El gatito activo controla el movimiento del carrusel; el gatito pasivo recibe un flujo visual aproximadamente equivalente, pero sin la misma relación causal de acción.

\lecturefigure{slide-01-two-kittens.jpg}{Experimento de los dos gatitos de Held y Hein: actuar activamente frente a observar pasivamente}{Diapositiva V001 recuperada de la grabación oficial; 00:03:18}

\begin{knowledgebox}{Lectura de la figura: no compares solo «cuánto se vio»}
En la figura, los dos gatitos reciben estímulos visuales similares, pero el movimiento del individuo pasivo de la izquierda lo impulsa el individuo activo de la derecha. Primero compara \textbf{quién decide la acción} y después el flujo visual; la diferencia decisiva no es la cantidad de píxeles, sino si existe entre action y observation una correspondencia causal que pueda aprenderse. El experimento respalda la importancia de la sensorimotor experience activa, pero por sí solo no demuestra que todas las capacidades cognitivas deban adquirirse desde cero mediante physical interaction.
\end{knowledgebox}

El ciclo cerrado mínimo entre agente y entorno puede escribirse así:
\[
o_t = \mathcal{O}(s_t),\qquad a_t \sim \pi_\theta(\cdot\mid h_t),\qquad
s_{t+1}\sim P(\cdot\mid s_t,a_t),\qquad h_{t+1}=f(h_t,o_{t+1},r_t).
\]
Donde $s_t$ es el estado real del entorno, $o_t$ es la observation visible, $a_t$ es la action, $\pi_\theta$ es la policy con parámetros $\theta$, $P$ es la dinámica del entorno, $r_t$ es la retroalimentación o reward, y $h_t$ es el estado interno construido a partir del historial. Observar pasivamente solo proporciona una parte de $o_t$; el aprendizaje activo dispone además de la cadena causal «qué $a_t$ ejecuté y por qué obtuve $o_{t+1}$ y $r_t$».

\teachervoice{El ponente llama a ChatGPT un «passive kitten» extremadamente potente: absorbió de internet una cantidad enorme de conocimiento humano, pero aprendió principalmente en calidad de observador. La propuesta de la clase no es abandonar el pretraining y dejar que el robot explore desde cero, sino añadir un ciclo cerrado activo sobre un conocimiento previo pasivo potente, para que el modelo aprenda las consecuencias de sus acciones, las posibilidades de acción y las tareas de largo plazo.}

\begin{importantbox}{El punto de partida de esta clase no es elegir entre pretraining y embodiment}
El preentrenamiento con datos de internet resuelve «qué sabe ya la humanidad»; la interacción activa resuelve «qué provocará mi acción con este cuerpo y en este entorno». Un generalist agent necesita combinar ambos: usar el foundation model para aportar el broad prior y la retroalimentación del entorno para calibrar el grounding, las skills y el control.
\end{importantbox}

\subsection{Resumen de la sección}

Esta sección fijó los límites de las fuentes, del tiempo y de la evidencia, y con el experimento de los dos gatitos estableció la tesis de toda la clase: la misma observation no equivale a la misma learning signal; el control de la acción y el ciclo cerrado de retroalimentación son indispensables. Todos los sistemas que siguen pueden entenderse como respuestas a una misma pregunta: cómo convertir el passive internet knowledge en una active competence ejecutable, verificable y acumulable.
\section{Qué es un generalist agent: objetivos, conocimiento y alcance de tareas}

La sección anterior explicó por qué hace falta un ciclo cerrado activo; esta sección define el alcance de las capacidades a las que debe servir ese ciclo. Un \term{generalist agent (agente generalista)} no es un specialist que «cubre muchos niveles en un benchmark», sino un sistema capaz de enfrentar objetivos nuevos que surgen continuamente, recurrir a un amplio conocimiento del mundo y transferir lo aprendido entre una gran cantidad de tareas heterogéneas. Aquí, \term{open-ended objective (objetivo abierto)} significa que el conjunto de tareas no tiene un punto final fijo, que los objetivos pueden ser planteados de forma continua por el usuario, el entorno o el propio agente, y que el criterio de éxito no necesariamente puede describirse por completo con una función booleana escrita a mano.

\subsection{De una reward fija a objetivos abiertos}

Un agente clásico de videojuegos suele recibir un entorno, un espacio de acciones y una reward fija; el entrenamiento puede ser muy potente, pero sigue optimizando solo un objetivo definido de antemano. Un mundo abierto exige que el sistema maneje objetos nuevos, combinaciones nuevas, restricciones nuevas y descripciones nuevas en lenguaje natural, por lo que no puede equipararse directamente «jugar muy bien» con generality. El curso usa tres requisitos para convertir la discusión, de un vocabulario publicitario, en una especificación de diseño verificable.

\lecturefigure{slide-02-generalist-agent-definition.jpg}{Los tres requisitos de capacidad de un generalist agent}{Diapositiva V002 recuperada del video oficial; 00:05:39}

\begin{knowledgebox}{Lectura de la figura: los tres requisitos deben cumplirse a la vez}
Primero observe a la izquierda open-ended objectives: las tareas no pueden agotarse en una lista fija; luego, en el centro, world knowledge: el sistema debe conocer objetos, reglas, sentido común y estrategias humanas; por último, massively multi-task: las capacidades deben compartirse entre una gran cantidad de tareas. Un benchmark multitarea que solo cumple el tercer requisito puede seguir siendo un conjunto cerrado, y un modelo conversacional que solo cumple el segundo carece de acción y de retroalimentación.
\end{knowledgebox}

Si las tareas provienen de una distribución $\mathcal{T}$ y cada tarea $\tau$ tiene un objetivo $g_\tau$ y una recompensa $R_\tau$, el objetivo de entrenamiento de una generalist policy puede abstraerse así:
\[
\max_\theta\;\mathbb{E}_{\tau\sim\mathcal{T}}\left[\mathbb{E}_{\pi_\theta}\left[\sum_{t=0}^{H_\tau}\gamma^t r_t^{(\tau)}\right]\right],
\qquad \text{y se exige mantener la capacidad de transferencia para }\tau'\notin\mathcal{T}_{\text{train}}\text{.}
\]
Donde $\tau$ denota la tarea, $g_\tau$ es un objetivo en lenguaje o multimodal, $H_\tau$ es la duración de la tarea, $r_t^{(\tau)}$ es la reward asociada a la tarea, $\gamma$ es el descuento temporal y $\theta$ son los parámetros compartidos de la policy. La dificultad de la fórmula no está en calcular la esperanza, sino en que $\mathcal{T}$ misma crece, la reward no siempre puede escribirse y las tareas de prueba no pueden resolverse solo memorizando la lista de entrenamiento.

\lecturefigure{slide-03-generalist-agent-requirements.jpg}{Los tres ingredients para construir un generalist agent}{Diapositiva V003 recuperada del video oficial; 00:05:45}

\begin{importantbox}{Los tres ingredients son relaciones de dependencia del sistema}
El entorno abierto determina «qué otras preguntas pueden plantearse», la internet-scale knowledge base determina «qué experiencia humana puede aprovecharse» y el foundation model determina «cómo representar de forma unificada la observation, el lenguaje, la memoria y la action». Si falta cualquiera de los tres eslabones, el sistema tiende a degradarse en un benchmark cerrado, en una exploración sin conocimiento o en un controlador especializado que no se transfiere.
\end{importantbox}

\teachervoice{El ponente primero compara specialist agents como los de Atari o Go y después presenta los tres requisitos de un generalist; con ello advierte al lector que no sustituya la definición de capacidad por una sola puntuación alta. Aunque un agent optimice un objetivo fijo hasta un nivel sobrehumano, mientras el objetivo, el entorno y la función de retroalimentación permanezcan congelados, sigue sin responder a la pregunta del mundo abierto: «¿qué es lo siguiente que vale la pena hacer?».}

\subsection{Ingredient 1: la complejidad del entorno es el límite de la capacidad}

El primer ingredient es el \term{open-ended environment (entorno abierto)}: objetos, terrenos, combinaciones de reglas y rutas de tareas lo bastante ricos para que el agente no agote pronto toda la estructura que puede descubrirse. El entorno no es un contenedor de pruebas pasivo; determina qué observations pueden producirse, qué actions pueden ejecutarse y qué curricula pueden construirse, y también determina si la evaluación es capaz de exponer la planificación a largo plazo y la generalización composicional.

\lecturefigure{slide-04-open-ended-environment.jpg}{El entorno abierto determina el límite superior de las capacidades que el agente puede aprender}{Diapositiva V004 recuperada del video oficial; 00:06:18}

\begin{knowledgebox}{Lectura de la figura: el entorno no es mejor por ser más grande, sino por tener estructura composicional}
La figura subraya que la agent capability está acotada superiormente por la environment complexity. Primero pregunte si el entorno incluye interacción entre objetos, dependencias de recursos, cambios de estado a largo plazo y varias formas de resolución; después, por el tamaño del mapa. Un espacio grande puramente aleatorio puede limitarse a aumentar el costo de búsqueda; lo que tiene verdadero valor formativo es que las reglas y los objetos puedan combinarse, de modo que las habilidades previas se conviertan en building blocks de tareas nuevas.
\end{knowledgebox}

\lecturefigure{slide-07-minecraft-open-world.jpg}{Minecraft como mundo abierto generado por procedimientos y composicional}{Diapositiva V007 recuperada del video oficial; 00:09:39}

La ventaja de Minecraft no es su «popularidad» en sí, sino que cuenta con procedural generation, dependencias de recolección y fabricación, combate, supervivencia, construcción y objetivos creativos casi arbitrarios. Al leer la figura, observe primero los recursos, las herramientas y la estructura espacial del mundo, y luego la vista en primera persona: una misma acción modifica el inventory, el terreno, los peligros y las tareas alcanzables después. Sigue siendo un simulator y no incluye toda la dinámica de un robot real, pero ofrece una capa intermedia escalable para la agent research de largo plazo y composicional.

\teachervoice{La afirmación del ponente es directa: el límite superior de la agent intelligence está restringido por la environment complexity. Este juicio también nos recuerda que, si un benchmark solo permite un horizon corto, una sola reward y un número limitado de objetos, por grande que sea el modelo, quizá solo esté memorizando con más rapidez tareas cerradas, en lugar de desarrollar una capacidad de exploración abierta.}

\subsection{Ingredient 2: Internet-scale World Knowledge}

El segundo ingredient es el \term{world knowledge (conocimiento del mundo)}: conocimiento previo reutilizable sobre propiedades de los objetos, reglas, recetas, relaciones causales, estrategias humanas y convenciones del lenguaje. En Minecraft, no se trata solo de datos enciclopédicos, sino también de «qué suele hacer un jugador ante esta situación», de las secuencias de operaciones en los videos, de las experiencias de fracaso en los foros y del conocimiento procedimental para descomponer un objetivo en lenguaje natural en pasos del árbol tecnológico.

\lecturefigure{slide-05-internet-knowledge-base.jpg}{La base de conocimiento de internet proporciona al agente un manual de referencia humano}{Diapositiva V005 recuperada del video oficial; 00:06:57}

\begin{knowledgebox}{Lectura de la figura: una knowledge base no equivale a una cantidad de páginas web}
El ícono de base de datos representa conocimiento que puede recuperarse, alinearse y transformarse. Primero distinga entre declarative knowledge (qué es un objeto) y procedural knowledge (cómo completar una tarea); después verifique si puede alinearse con la observation, el inventory y el goal actuales del agent. Un texto web sin grounding puede ser correcto y, aun así, no convertirse directamente en una action en el mundo actual.
\end{knowledgebox}

\subsection{Ingredient 3: el Foundation Model como interfaz unificada}

El tercer ingredient es un flexible foundation model. Su función aquí no es solo generar lenguaje natural, sino conectar objetivos, estados visuales, código, retroalimentación, memoria y action representation en una misma interfaz de modelado de secuencias. El modelo debe poder recurrir a un broad prior y, a la vez, aceptar la executable evidence que devuelve el entorno; de lo contrario, sigue siendo solo un sistema de preguntas y respuestas fuera de línea.

\lecturefigure{slide-06-foundation-model-for-agents.jpg}{El foundation model conecta los objetivos abiertos con la retroalimentación del entorno}{Diapositiva V006 recuperada del video oficial; 00:07:36}

\begin{warningbox}{Un «foundation model» no trae grounding automáticamente}
La escala de parámetros y el web pretraining pueden mejorar los conocimientos previos de lenguaje y visión, pero si el objeto actual puede manipularse, si el código puede ejecutarse, si la acción es segura y si se están aprovechando fallas de la reward, todo ello debe verificarse mediante la environment interface y la retroalimentación. Tomar el contenido generado directamente como una capacidad real es el category error más común en los sistemas de agents.
\end{warningbox}

\subsection{Resumen de la sección}

La definición de generalist agent comprende objetivos abiertos, conocimiento del mundo y capacidad multitarea a gran escala; construirlo requiere un entorno abierto, una base de conocimiento de internet y un flexible foundation model. Los tres determinan, respectivamente, el espacio de problemas, la fuente del conocimiento previo y la interfaz unificada; los trabajos posteriores, MineDojo, MineCLIP, Voyager, Eureka y VIMA, completan cada uno distintos eslabones de esta recipe.
\section{MineDojo: convertir Minecraft en una Agent Research Platform abierta}

Con la recipe definida, el curso usa primero MineDojo para llevar los requisitos abstractos a una infraestructura concreta. MineDojo ofrece a la vez un environment API, un ecosistema de tareas y una internet knowledge database. Así, quien investiga no solo puede entrenar una policy, sino también estudiar la generación de tareas, el reward learning, la recuperación de conocimiento y las habilidades de largo plazo. Para entender esta sección hay que distinguir tres capas: qué puede simular el entorno, cómo se definen las tareas y qué tipo de supervisión pueden aportar los materiales de internet.

\subsection{Entorno, agente y base de conocimiento: un conjunto de tres piezas}

El diagrama del sistema de MineDojo divide el mundo abierto en tres objetos que se restringen entre sí: el Minecraft environment produce la información visual y el estado, el generalist agent elige la action y la internet knowledge base aporta demostraciones humanas y conocimiento previo en texto. No basta con entregarle páginas web al agent: hay que construir interfaces entre goal, observation, knowledge y control.

\lecturefigure{slide-08-minedojo-three-components.jpg}{Entorno abierto, agente y base de conocimiento de internet en MineDojo}{Diapositiva V008 recuperada de la grabación oficial; 00:10:06}

\begin{knowledgebox}{Lectura de la figura: seguir el flujo de datos, no los íconos}
Desde el entorno, el agente recibe la imagen en primera persona, el inventory y el estado; desde la base de conocimiento, obtiene materiales humanos como videos, Wiki y Reddit; a su vez, la action del agent modifica el entorno. La pregunta central es cómo se alinean ambos tipos de entrada con un mismo objetivo, no si los tres componentes son «todos muy grandes».
\end{knowledgebox}

\lecturefigure{slide-09-three-thousand-tasks-title.jpg}{MineDojo organiza la investigación de mundo abierto con más de 3,000 tareas}{Diapositiva V009 recuperada de la grabación oficial; 00:10:24}

«3,000+ tasks» es el punto de entrada a las tareas; no debe leerse como si la cantidad ya demostrara la generality. El valor de la escala de tareas está en cubrir distintos recursos, lugares, objetos, horizontes temporales y condiciones de éxito, de modo que se puedan observar la transfer y la compositionality. Si todas las tareas siguieran compartiendo un mismo patrón superficial, la cantidad por sí sola no produciría capacidad de mundo abierto.

\subsection{Programmatic Tasks y Creative Tasks}

El espacio de tareas de MineDojo contiene dos tipos de objetivos complementarios. Una programmatic task puede evaluarse automáticamente a partir del simulator state, por ejemplo, si el inventory contiene cierto objeto o si se derrotó a cierta criatura. Una creative task, en cambio, puede pedir «construir una casa bonita» o «apagar un incendio de forma creativa»; su criterio de éxito tiene componentes semánticos y estéticos, y no puede escribirse simplemente como un predicado de estado.

\lecturefigure{slide-10-task-categories.jpg}{Simulator, tareas programáticas y tareas creativas en MineDojo}{Diapositiva V010 recuperada de la grabación oficial; 00:10:48}

\begin{knowledgebox}{Términos clave: responsabilidades de los tres tipos de infraestructura}
\begin{tabularx}{\textwidth}{@{}L{0.20\textwidth}X X@{}}
\toprule
Componente & Problema que resuelve & Relación con el hilo del curso \\
\midrule
Versatile simulator & Un mundo observable, controlable y paralelizable & Aporta el ciclo cerrado activo y experimentos reproducibles \\
Programmatic task & Objetivos explícitos cuyo estado puede verificarse & Adecuada para reward automática y entrenamiento a gran escala \\
Creative task & Objetivos abiertos para los que es difícil escribir a mano un success predicate & Obliga al sistema a aprender una reward semántica y las preferencias humanas \\
\bottomrule
\end{tabularx}
\end{knowledgebox}

\lecturefigure{slide-11-inventory-tech-tree-combat.jpg}{Inventory, technology tree y combat forman un espacio de tareas combinatorio}{Diapositiva V011 recuperada de la grabación oficial; 00:11:18}

Esta diapositiva debe leerse primero por sus dependencias: obtener objetos avanzados suele requerir recolectar materias primas, fabricar herramientas, desbloquear recetas y enfrentar peligros; por eso, las tareas no son labels independientes, sino que comparten un mismo árbol tecnológico. Los objetos, herramientas y criaturas de la figura son solo ejemplos; lo que de verdad sustenta la transferencia es que las skills previas puedan combinarse en un nuevo plan de horizon más largo.

\lecturefigure{slide-12-creative-task-examples.jpg}{Las tareas creativas carecen de un criterio automático de éxito sencillo}{Diapositiva V012 recuperada de la grabación oficial; 00:11:57}

\begin{warningbox}{Creative no significa «asignar la reward de cualquier manera»}
Las tareas creativas son más difíciles de evaluar automáticamente, pero siguen requiriendo evidence explícita: el texto del objetivo, la semántica del video, las preferencias humanas, las restricciones del entorno o un learned reward model. Si se usa solo una puntuación ambigua, la policy puede encontrar un shortcut ajeno a la intención humana. Una evaluación abierta debe informar al mismo tiempo el origen de la reward, los casos de fallo y la revisión humana.
\end{warningbox}

\lecturefigure{slide-13-firefight-open-ended-demo.jpg}{La tarea abierta de apagar un incendio muestra el objetivo semántico y la secuencia de acciones}{Diapositiva V013 recuperada de la grabación oficial; 00:12:57}

Al ver la demo, primero hay que distinguir entre «la imagen parece razonable» y «la tarea se completó de forma verificable». El video muestra al agent frente al fuego, la estructura del entorno y la elección de herramientas, pero un único éxito no ofrece tasa de éxito, alcance de generalización ni reward robustness. Sobre todo ilustra que un objetivo abierto requiere conectar la semántica visual, el inventory y las acciones de varios pasos, en lugar de solo predecir el siguiente cuadro.

\teachervoice{El ponente distingue expresamente entre programmatic y creative tasks: en las primeras es fácil escribir el success criterion; en las segundas, incluso decidir «qué cuenta como completado» puede requerir un juicio semántico. Esta distinción explica directamente por qué después se introduce MineCLIP: la learned reward no es un módulo adicional, sino la interfaz por la que los objetivos abiertos entran en el RL.}

\subsection{Base de datos de internet: de la tabla de objetos a las preguntas humanas}

El entorno aporta experiencia, pero aprender desde cero, por trial-and-error, todas las recetas y estrategias resulta extremadamente costoso. Por eso MineDojo trata YouTube, Wiki y Reddit como un acervo de experiencia humana: los videos aportan conducta e información visual, Wiki aporta hechos estructurados y los foros aportan fallas, estrategias y problemas de cola larga. Las tres fuentes difieren en ruido y en dificultad de grounding, por lo que no pueden concatenarse simplemente como tokens homogéneos.

\lecturefigure{slide-14-internet-knowledge-sources.jpg}{YouTube, Wiki y Reddit constituyen fuentes de conocimiento complementarias}{Diapositiva V014 recuperada de la grabación oficial; 00:13:06}

\begin{knowledgebox}{Lectura de la figura: qué aporta cada una de las tres fuentes}
YouTube destaca en demostraciones dinámicas, pero no tiene action labels confiables; Wiki destaca en objetos, recetas y reglas, pero rara vez muestra la recuperación ante fallos; Reddit destaca en la resolución de problemas y la experiencia de cola larga, pero contiene ruido y juicios subjetivos. Un generalist agent necesita source-aware retrieval, en lugar de tratar todo el contenido como un mismo tipo de supervisión.
\end{knowledgebox}

\lecturefigure{slide-15-internet-knowledge-scale.jpg}{Materiales de internet a gran escala cubren el espacio visual y semántico de Minecraft}{Diapositiva V015 recuperada de la grabación oficial; 00:13:18}

Las miniaturas densas de esta diapositiva ilustran ante todo la coverage, no labels limpios. Al leer la figura, compare la variedad de objetos, escenas y puntos de vista; esto respalda que «internet contiene abundante conocimiento previo», pero no demuestra que los materiales ya estén alineados con el agent state, y mucho menos que todas las conductas de los videos sean correctas o reproducibles. La escala debe auditarse junto con la eliminación de duplicados, la calidad, los derechos de autor y la pertinencia para la tarea.

\lecturefigure{slide-16-minecraft-item-catalog.jpg}{Los catálogos de objetos, criaturas y acciones aportan conocimiento a nivel de entidad}{Diapositiva V016 recuperada de la grabación oficial; 00:14:15}

El catálogo de objetos ayuda al modelo a construir un entity vocabulary: qué objetos pueden recolectarse, combinarse, equiparse o resultan peligrosos. Primero conviene ver los límites entre categorías y después las posibilidades de acción que ofrece cada objeto respecto a cada action. Conocer el nombre no equivale a saber qué píxeles del punto de vista actual corresponden al objeto, ni a conocer las precondiciones de la acción; por eso siguen siendo necesarios el perception grounding y el environment feedback.

\lecturefigure{slide-17-crafting-knowledge.jpg}{El crafting knowledge conecta el objeto objetivo con sus materias primas y pasos}{Diapositiva V017 recuperada de la grabación oficial; 00:14:24}

Las recetas de fabricación convierten un declarative fact en una procedural constraint. Al leer la figura, rastree hacia atrás, desde el objeto objetivo, las materias primas, las herramientas y las condiciones de la mesa de trabajo; esta dependencia puede usarse para la plan decomposition y también para el curriculum. El diagrama respalda que «internet puede aportar pasos estructurados», pero durante la ejecución todavía hay que verificar el inventory, la ubicación y la recuperación ante fallos.

\lecturefigure{slide-18-potion-recipe-graph.jpg}{El grafo de recetas de pociones muestra cadenas largas de dependencias y planes combinados}{Diapositiva V018 recuperada de la grabación oficial; 00:14:33}

La dificultad de este tipo de grafo de dependencias no está en memorizar una recipe aislada, sino en mantener el estado a lo largo de varios objetos intermedios. Primero se identifica el nodo objetivo y luego se recorren las aristas hacia atrás para obtener cada ingredient necesario; después, el orden lógico del grafo se traduce en actions de navigation, collection y crafting dentro del mundo. Cualquier recurso faltante puede obligar al agent a replanificar.

\lecturefigure{slide-19-wiki-reddit-database.jpg}{Las páginas de Wiki y las discusiones de Reddit aportan reglas y experiencia de cola larga}{Diapositiva V019 recuperada de la grabación oficial; 00:14:42}

Las páginas web y los hilos de discusión contienen explicaciones humanas, alternativas y experiencias de fallos; son una teacher signal importante además de los videos. Al leer la figura, primero separe las páginas de hechos de las publicaciones de opinión y después revise la fecha de publicación, la versión y el contexto. Las respuestas del foro pueden inspirar estrategias, pero no deben ejecutarse sin verificación; el environment result es la evidencia final en la tarea actual.

\lecturefigure{slide-20-human-open-ended-questions.jpg}{Las preguntas abiertas planteadas por personas amplían la distribución de tareas}{Diapositiva V020 recuperada de la grabación oficial; 00:15:09}

Estas preguntas muestran que la distribución de objetivos proviene de jugadores reales y no de un conjunto limitado de plantillas diseñadas de antemano por los investigadores. Primero vea si la tarea requiere sentido común, planificación o creatividad, y después si puede calificarse automáticamente. Las preguntas sirven para construir benchmarks de task proposal y de retrieval, pero también exponen el problema de la evaluación: los objetivos en lenguaje pueden ser ambiguos y el éxito suele admitir varias realizaciones equivalentes.

\teachervoice{El ponente llama a los materiales de internet un reference manual: cuando una persona no sabe cómo hacer algo en Minecraft, también busca en Wiki, en videos y en foros. Lo esencial para el agent no es «memorizar todo internet», sino encontrar el conocimiento pertinente para un goal y un state concretos, y juzgar con los resultados de la ejecución si esa sugerencia realmente aplica.}

\lecturefigure{slide-21-generalist-agent-transition.jpg}{Del entorno y la base de conocimiento a un generalist agent que puede aprender}{Diapositiva V021 recuperada de la grabación oficial; 00:15:18}

Esta diapositiva de transición reduce el problema a lo siguiente: si ya se tienen un entorno abierto y conocimiento humano, ¿cómo se entrena un agent capaz de entender objetivos en lenguaje natural? La respuesta, en primer lugar, no es predecir acciones directamente, sino aprender un reward model capaz de alinear videos con objetivos en texto; ese modelo convierte objetivos semánticos difíciles de escribir a mano en retroalimentación que el RL puede usar.

\subsection{Resumen de la sección}

MineDojo reúne en una misma plataforma el entorno abierto, el ecosistema de tareas y el conocimiento de internet. Las programmatic tasks permiten el entrenamiento automático, mientras que las creative tasks revelan la carencia de evaluación semántica; YouTube, Wiki y Reddit aportan conocimiento complementario, pero este debe alinearse con el estado actual y verificarse mediante la ejecución. Esto establece la base de datos e interfaces para la learned reward de MineCLIP y para la línea de code/memory de Voyager.
\section{MineCLIP: convertir el Internet Video en una Language-conditioned Reward}

La carencia más difícil de la sección anterior es que el creative goal no tiene una reward fácil de escribir a mano. La estrategia de MineCLIP es aprender una representación común de fragmentos de video y descripciones textuales: si el contenido del video corresponde mejor al texto objetivo, la similitud es mayor, y puede usarse como dense feedback para el aprendizaje de la policy. Un \term{reward model (modelo de recompensa)} es un modelo que asigna a una trajectory, una observation o un outcome una retroalimentación escalar; aproxima el objetivo humano, pero no es el objetivo mismo.

\subsection{Contrastive Learning e InfoNCE}

El \term{contrastive learning (aprendizaje contrastivo)} aprende representaciones acercando las muestras que corresponden entre sí y alejando las que no corresponden. MineCLIP codifica los fragmentos de video de Minecraft como $v_i$ y el texto correspondiente como $t_i$; durante el entrenamiento, el mismo par es positive y los demás textos o videos del lote funcionan como negative, de modo que los video--text pair con semántica coherente quedan más cerca en el embedding space.

\lecturefigure{slide-22-mineclip-contrastive-reward.jpg}{MineCLIP aprende la alineación video--lenguaje y produce una similitud semántica}{Diapositiva V022 recuperada de la grabación oficial; 00:16:51}

\begin{knowledgebox}{Lectura de la figura: distinguir primero la señal de entrenamiento de la reward en despliegue}
En la etapa de entrenamiento se usan pares de video clip y transcript, de modo que los pares que corresponden obtengan puntajes altos; en la etapa de despliegue, el segmento más reciente de observation del agent y el texto objetivo se envían al modelo para obtener la similarity reward. En la figura, un puntaje alto indica mayor cercanía semántica, no que la tarea esté necesariamente completada; el modelo puede preferir atajos visuales, el fondo o coocurrencias frecuentes, por lo que se requieren un outcome check y una auditoría humana.
\end{knowledgebox}

Una InfoNCE loss video-to-text simplificada es:
\[
\mathcal{L}_{\mathrm{NCE}}
=-\frac{1}{B}\sum_{i=1}^{B}
\log\frac{\exp(\operatorname{sim}(v_i,t_i)/\tau)}
{\sum_{j=1}^{B}\exp(\operatorname{sim}(v_i,t_j)/\tau)}.
\]
Donde $B$ es el tamaño del lote, $v_i$ es el embedding del $i$-ésimo video, $t_i$ es el embedding del texto correspondiente, $\operatorname{sim}$ suele ser la cosine similarity de vectores normalizados y $\tau$ es la temperature. La loss solo exige que los pares correspondientes estén relativamente más cerca; no garantiza que la similarity esté calibrada como una probabilidad real de éxito.

Para un texto objetivo $g$, la reward semántica que usa la policy puede escribirse así:
\[
r_t^{\text{MineCLIP}}=\operatorname{sim}\!\left(E_v(o_{t-k:t}),E_t(g)\right),
\qquad
J(\theta)=\mathbb{E}_{\pi_\theta}\left[\sum_t\gamma^t r_t^{\text{MineCLIP}}\right].
\]
Donde $E_v$ y $E_t$ son, respectivamente, los encoder de video y de texto, $o_{t-k:t}$ es la observation window de los $k$ pasos más recientes, $g$ es el objetivo en lenguaje, $r_t^{\text{MineCLIP}}$ es la learned reward y $J(\theta)$ es el policy objective. Cuanto más dense es la reward, más fácil resulta la exploración; pero la policy también amplifica activamente el reward bias.

\subsection{Entrenar una Language-conditioned Policy con una Learned Reward}

MineCLIP por sí solo no es un controlador completo. El sistema sigue necesitando una RL policy que elija la action a partir de la observation actual y que actualice sus parámetros con la learned reward. Por ello, el objetivo en lenguaje interviene en dos lugares: en uno define la dirección semántica de la reward; en el otro puede servir como policy condition, para que una misma red ejecute tareas distintas según la instruction.

\lecturefigure{slide-23-rl-with-mineclip.jpg}{La reward de MineCLIP impulsa el language-conditioned reinforcement learning}{Diapositiva V023 recuperada de la grabación oficial; 00:17:03}

\begin{lstlisting}[caption={Pseudocódigo conceptual del ciclo de entrenamiento de RL condicionado por MineCLIP}]
for goal in sampled_language_goals:
    obs = env.reset(goal)
    recent_frames = []
    while not env.done():
        action = policy.sample(obs, goal)
        next_obs = env.step(action)
        recent_frames.append(next_obs.frame)
        reward = mineclip.similarity(recent_frames[-k:], goal)
        replay.add(obs, action, reward, next_obs, goal)
        policy.update(replay)
        obs = next_obs
\end{lstlisting}

\begin{warningbox}{El reward hacking se agrava a medida que la policy se vuelve más capaz}
Si MineCLIP confunde cierto fondo, un movimiento de cámara o un objeto local con el éxito, el RL buscará ese shortcut de forma sistemática. Es indispensable registrar a la vez el programmatic state, la revisión humana de los videos, distintas seed, objetivos adversariales y la reward trajectory; «la reward sube» solo indica que la señal de entrenamiento mejoró, no equivale automáticamente a que se haya cumplido la intención humana.
\end{warningbox}

\teachervoice{El ponente describe MineCLIP como el puente entre el passive internet video y el active reinforcement learning: el video aporta la similitud semántica y luego el agent aprende a controlar mediante la interacción con el entorno. Este puente es importante, pero en la clase tampoco se afirma que la learned reward sea perfecta; sigue limitada por sesgos de los datos, ambigüedad de los objetivos y shortcut visuales.}

\subsection{Resultados: éxito en las tareas y robustez visual}

La diapositiva de resultados debe responder dos preguntas distintas: si la reward de MineCLIP basta para entrenar una policy que cumpla varios objetivos en lenguaje, y si la policy solo memoriza los escenarios de entrenamiento. Para la primera, la evidencia es el success de las tareas o la comparación con un baseline; para la segunda, la evidencia es la zero-shot visual condition. Ambas deben acotarse a las tareas evaluadas, las seed y la Minecraft distribution.

\lecturefigure{slide-24-mineclip-results.jpg}{Resultados del RL condicionado por MineCLIP en varios tipos de tareas}{Diapositiva V024 recuperada de la grabación oficial; 00:19:00}

\begin{knowledgebox}{Lectura de la figura: revisar primero la definición de las columnas y después la mejora relativa}
La tabla compara el indicador de success de distintas tareas y métodos. Primero hay que confirmar que la tarea de cada columna use el mismo evaluation protocol y luego comparar la reward de MineCLIP con las demás reward/baseline; la conclusión principal es que la representación video--lenguaje puede aportar una señal de entrenamiento útil, no que todas las tareas estén resueltas. El promedio oculta las tareas fallidas y la seed variance, y no permite inferir generality en un mundo abierto.
\end{knowledgebox}

\lecturefigure{slide-25-policy-robustness.jpg}{La policy aprendida conserva cierta robustez en condiciones visuales no vistas}{Diapositiva V025 recuperada de la grabación oficial; 00:19:57}

\begin{knowledgebox}{Lectura de la figura: la robustez proviene de la invariancia de la tarea, no de una capacidad OOD arbitraria}
La figura muestra distintos climas, iluminaciones, biomas o condiciones visuales. Primero conviene comparar si la conducta objetivo se mantiene y después verificar si esos cambios siguen perteneciendo a las reglas del mundo de entrenamiento. El resultado respalda que la policy no quedó totalmente ligada a una sola imagen, pero no demuestra que pueda transferirse a un juego nuevo, a un nuevo espacio de acciones ni a un robot real; la visual robustness y la embodiment transfer son niveles distintos.
\end{knowledgebox}

\teachervoice{La postura de la clase ante los resultados es que «la reward representation es útil», no que «Minecraft ya está resuelto». El ponente destaca la robustez bajo unseen visual condition porque se acerca más a la foundation-model transfer; pero también califica MineDojo, Voyager y Eureka como apenas un primer paso, scratching the surface.}

\subsection{Resumen de la sección}

MineCLIP usa el contrastive video--language learning para convertir objetivos en lenguaje natural en una dense reward, con la que después una RL policy aprende a controlar en el entorno. Demuestra que el internet video puede ser una fuente de supervisión para el active learning y, al mismo tiempo, expone el riesgo central de la learned reward: la similitud no equivale al éxito, la policy amplifica los sesgos del reward model y los resultados deben interpretarse junto con el estado programático y la auditoría humana.
\section{Voyager: Code-as-Action, Memory y Automatic Curriculum}

MineCLIP conecta la semántica con la policy de bajo nivel mediante la reward, pero la exploración abierta de horizon largo todavía requiere una action representation más componible. Voyager opta por \term{code-as-action (el código como acción)}: el LLM no emite instrucciones de teclado cuadro por cuadro, sino que genera programas en JavaScript que llaman a la API de Minecraft. Los programas pueden incluir bucles, condicionales, manejo de errores y llamadas a skills existentes, con lo que comprimen secuencias largas de acciones en unidades simbólicas ejecutables y verificables.

\subsection{Por qué aplicar «stringify» al video y al mundo}

Los videos de internet contienen comportamientos abundantes, pero carecen de action labels limpias; el problema inverso de pasar de píxeles a control continuo suele tener múltiples soluciones. El curso usa «stringify» para expresar una estrategia de interfaz: convertir el estado del entorno, el inventory, los errores y las API disponibles en contexto de lenguaje o código, para que el LLM razone en el token space, donde es más fuerte, y después las herramientas ejecuten y devuelvan resultados reales.

\lecturefigure{slide-26-stringify-video.jpg}{De trayectorias visuales a una interfaz simbólica que un modelo de lenguaje puede procesar}{Diapositiva V026 recuperada de la grabación oficial; 00:21:39}

\begin{knowledgebox}{Lectura de la figura: la stringification es ingeniería de interfaces, no una conversión de información sin pérdida}
La figura pasa del video dinámico a módulos de texto o simbólicos. La conversión debe elegir qué state se expone al modelo, cómo se describen las relaciones espaciales y qué action se encapsulan como API. Mejora la legibilidad y la componibilidad, pero puede perder información fina de geometría, temporalidad y contacto; por ello se adapta mejor a skills de Minecraft de alto nivel y no debe considerarse directamente una respuesta general para todo el control de bajo nivel de robots.
\end{knowledgebox}

\lecturefigure{slide-27-mineflayer-javascript.jpg}{La Mineflayer JavaScript API proporciona una interfaz de control ejecutable para Minecraft}{Diapositiva V027 recuperada de la grabación oficial; 00:22:15}

Mineflayer encapsula la navegación, la recolección, la fabricación y la interacción como llamadas de programa. Al leer la figura, primero se observan los parámetros de las funciones y los nombres de los objetos, y después cómo se devuelven los resultados de ejecución; la API permite probar y reutilizar el código, y también establece una safety boundary. Al mismo tiempo, la calidad del encapsulamiento determina el techo de capacidad: los detalles que la API no expone no los puede recuperar el LLM solo con un prompt más ingenioso.

\teachervoice{El ponente señala que el video es muy rico en información, pero no tiene action labels; por eso Voyager no intenta recuperar directamente las operaciones de teclado de cada cuadro, sino que aprovecha la symbolic interface de Mineflayer. Esta decisión conecta «entender el objetivo» con «ejecutar el programa», y además permite que el syntax error, el runtime error y el task progress se conviertan en retroalimentación legible para el modelo.}

\subsection{Voyager Data Flow e Iterative Prompting}

Voyager no termina después de que un solo prompt genera un fragmento de código. Mantiene el state actual, el objetivo, las skills disponibles y la retroalimentación de ejecución, y le pide repetidamente al actor que genere o revise programas; el environment, a su vez, restringe al modelo con el estado real, los errores y el progress. Aquí, el \term{critic (crítico)} es el módulo que determina si la tarea se completó, dónde falló y qué dirección de mejora seguir; no se limita a dar una opinión de estilo.

\lecturefigure{slide-28-voyager-data-flow.jpg}{Flujo de datos de Voyager entre curriculum, code action, environment y memory}{Diapositiva V028 recuperada de la grabación oficial; 00:23:12}

\begin{knowledgebox}{Lectura de la figura: seguir cuatro tipos de estado a lo largo del ciclo cerrado}
El curriculum produce una task; el actor lee la observation y las relevant skills y genera code; el environment ejecuta y devuelve el world state y los error; el critic evalúa el progress; las skills exitosas se escriben en la library. Lo importante no es cuántas llamadas a GPT-4 hay, sino si los cuatro tipos de estado (task, execution, evaluation y memory) forman un ciclo cerrado.
\end{knowledgebox}

\lecturefigure{slide-29-code-as-action.jpg}{GPT-4 genera programas ejecutables para controlar un embodied agent}{Diapositiva V029 recuperada de la grabación oficial; 00:23:27}

Esta diapositiva sustituye los action tokens tradicionales por programas completos. Primero se observan las llamadas a funciones, los bucles y las verificaciones en la salida del actor, y después el comportamiento de largo plazo del agent en el mundo; el código comprime cientos de pasos de acción en una sola skill, pero también introduce riesgos de software como prompt injection, bucles infinitos, llamadas destructivas y dependency errors, por lo que debe ejecutarse dentro de un sandbox y con una API allowlist.

\lecturefigure{slide-30-iterative-prompting.jpg}{La retroalimentación del entorno, los errores de ejecución y la autoverificación impulsan la iteración del código}{Diapositiva V030 recuperada de la grabación oficial; 00:24:09}

\begin{warningbox}{Iterative prompting no significa que los errores desaparezcan solos}
La retroalimentación solo es útil cuando es observable, atribuible y llega a tiempo. Si falta el environment state, el critic se equivoca o el programa altera el mundo de forma irrecuperable, el siguiente prompt puede seguir parchando sobre una premisa errónea. El sistema debe limitar el retry budget, conservar el execution trace y someter a aprobación por separado las action irreversibles.
\end{warningbox}

\begin{lstlisting}[caption={Ciclo conceptual de Voyager: generación de código, retroalimentación de ejecución y skill memory}]
task = curriculum.propose(agent_state, skill_library)
while retry_budget > 0:
    skills = skill_library.retrieve(task, agent_state)
    program = actor.write_code(task, agent_state, skills)
    result = sandbox.execute(program, allowed_api=mineflayer_api)
    verdict = critic.check(task, result, agent_state)
    if verdict.success:
        skill_library.store(task, program, result.preconditions)
        break
    agent_state = result.updated_state
    actor.reflect(verdict.feedback, result.errors)
    retry_budget -= 1
\end{lstlisting}

\teachervoice{En clase se describe Voyager como una no-gradient architecture: en tiempo de ejecución no se actualizan los parámetros del modelo, sino que el comportamiento mejora mediante el prompt, la retroalimentación de ejecución y la memory externa. Aquí el «learning» ocurre principalmente en la skill library y en el curriculum state, y no debe confundirse con un neural-network weight update.}

\subsection{Critic y Skill Library: convertir un éxito aislado en una capacidad reutilizable}

Si cada tarea partiera de un prompt en blanco, la exploración abierta pagaría repetidamente el mismo costo. La \term{skill library (biblioteca de habilidades)} es una executable memory persistente: guarda los programas exitosos, su propósito, sus precondiciones y sus descripciones para la recuperación, de modo que las tareas posteriores puedan combinar habilidades anteriores. El critic, por su parte, convierte la environment evidence en juicios de finalización y señales de revisión, y evita que un código se escriba en la memory solo porque «parece correcto».

\lecturefigure{slide-31-critic-self-reflection.jpg}{El critic genera self-reflection a partir del avance de la tarea}{Diapositiva V031 recuperada de la grabación oficial; 00:24:45}

\begin{knowledgebox}{Lectura de la figura: el critic debe citar estados verificables}
Un feedback eficaz debe indicar qué objeto falta, por qué la posición actual no cumple la condición y en qué paso falló la ejecución, en lugar de decir vagamente «try again». Primero se revisa qué environment observation cita el verdict y después si la sugerencia puede cambiar el siguiente programa. El propio critic también puede alucinar, por lo que conviene que el success lo determinen conjuntamente un state predicate y un juicio semántico.
\end{knowledgebox}

\lecturefigure{slide-32-skill-library-overview.jpg}{La skill library hace persistentes executable procedures cada vez más complejos}{Diapositiva V032 recuperada de la grabación oficial; 00:25:27}

El valor de la biblioteca de habilidades proviene de la compositionality: las habilidades tempranas de recolección, fabricación y navegación pueden ser invocadas por tareas de nivel superior. Al leer la figura, primero se revisa si la granularidad de las habilidades es lo bastante estable y después las dependencias entre skills; si son demasiado finas, el prompt se infla; si son demasiado gruesas, resulta difícil reutilizarlas. Cada skill necesita además precondiciones, side effects y failure modes, no solo un nombre de función.

Para una task query $q$, una skill description $d_i$ y el estado actual $x$, una regla de recuperación simplificada es:
\[
i^*=\arg\max_i\left[\operatorname{sim}(E(q,x),E(d_i))
-\lambda\,\operatorname{cost}(i,x)\right].
\]
Donde $E$ es el embedding encoder, $\operatorname{sim}$ mide la relevancia semántica entre la tarea y la habilidad, $\operatorname{cost}(i,x)$ representa el costo en recursos o el riesgo de invocar la habilidad en el estado actual, $\lambda$ controla el equilibrio entre relevancia y costo, e $i^*$ es la skill recuperada. Usar solo la similitud semántica puede recuperar programas de nombre parecido cuyas precondiciones no se cumplen.

\lecturefigure{slide-33-retrieve-old-skills.jpg}{Recuperar y combinar skills anteriores para completar tareas nuevas más complejas}{Diapositiva V033 recuperada de la grabación oficial; 00:26:36}

Esta diapositiva muestra la ruta retrieve--compose: una tarea nueva activa primero programas anteriores relevantes y luego el actor genera la lógica de composición. Primero se verifica si las habilidades anteriores realmente se invocan y después se comparan los costos en tokens, errores y tiempo entre generar desde cero y reutilizar. Esto respalda que la external memory puede acumular competence, pero no demuestra que los programas sigan siendo confiables cuando cambia la versión del entorno; la library necesita regression tests y mecanismos de invalidación.

\teachervoice{El ponente enfatiza que la skill library hace que Voyager no sea una «demo de ChatGPT de un solo uso». El código exitoso se guarda y se recupera en tareas más difíciles, de modo que el sistema logra un crecimiento de capacidad a largo plazo sin actualizar pesos. La idea más transferible aquí es almacenar la agent memory como un artifact ejecutable y verificable, en lugar de guardar solo resúmenes en lenguaje natural.}

\subsection{Automatic Curriculum: quién decide qué aprender a continuación}

Un mundo abierto no tiene un orden fijo de niveles, por lo que también se necesita un \term{automatic curriculum (currículo automático)}: a partir del inventory actual, las habilidades desbloqueadas, el historial de exploración y la novelty, propone la siguiente tarea «factible ahora, pero aún no completada». Controla a la vez la dificultad y la dirección de la exploración; si la tarea es demasiado difícil, se desperdician llamadas, y si es demasiado fácil, no se amplían las capacidades.

\lecturefigure{slide-34-automatic-curriculum.jpg}{El automatic curriculum amplía las capacidades a lo largo del árbol tecnológico de Minecraft}{Diapositiva V034 recuperada de la grabación oficial; 00:27:18}

\begin{knowledgebox}{Lectura de la figura: el curriculum es una task policy}
La figura va de los recursos básicos a los objetos avanzados y muestra que las tareas tienen dependencias previas. Primero se observan el inventory actual y las unlocked skills, y después si las tareas candidatas conducen a estados nuevos; un buen curriculum no persigue simplemente la reward más alta, sino que busca el equilibrio entre feasibility, novelty, usefulness y risk.
\end{knowledgebox}

\lecturefigure{slide-35-proposed-next-task.jpg}{El curriculum propone, según el estado, una siguiente tarea factible y novedosa}{Diapositiva V035 recuperada de la grabación oficial; 00:28:30}

La selección de tareas puede abstraerse así:
\[
\tau_{k+1}=\arg\max_{\tau\in\mathcal{C}(x_k)}
\bigl[\alpha F(\tau\mid x_k)+\beta N(\tau\mid M_k)+\eta U(\tau)-\rho R(\tau)\bigr].
\]
Donde $\mathcal{C}(x_k)$ es el conjunto de tareas candidatas en el estado actual $x_k$, $F$ es la feasibility, $N$ es la novelty respecto de la memory $M_k$, $U$ es el valor de reutilización futura, $R$ es el riesgo o el costo esperado, y $\alpha,\beta,\eta,\rho$ son pesos. El Voyager real usa una LLM proposal; la fórmula sirve para mostrar las dimensiones de evaluación y no significa que el sistema calcule explícitamente estos escalares.

\subsection{Resultados del ciclo cerrado: velocidad en el árbol tecnológico y cobertura del mapa}

Solo al ensamblar curriculum, actor, critic y skill library se convierte Voyager en un sistema de exploración abierta. Las diapositivas de resultados miden por separado el progress en el árbol tecnológico y la spatial exploration; el primero se aproxima al desbloqueo de habilidades y recursos, y la segunda a la cobertura del mapa. Ambos tipos de métricas son complementarios, pero los dos son solo proxies dentro de Minecraft.

\lecturefigure{slide-36-voyager-all-pieces.jpg}{El curriculum, el prompting y la skill library de Voyager forman un ciclo cerrado}{Diapositiva V036 recuperada de la grabación oficial; 00:29:03}

\begin{importantbox}{El ciclo cerrado mínimo de Voyager}
El curriculum propone tareas, el actor escribe código, el environment devuelve evidencia de ejecución, el critic evalúa el avance, los programas exitosos entran en la skill library y la nueva memory, a su vez, modifica la siguiente tarea y el siguiente fragmento de código. Si falta cualquier eslabón, el sistema degenera en exploración sin dirección, generación de un solo uso, ejecución no verificable o una demo incapaz de acumular.
\end{importantbox}

\lecturefigure{slide-37-open-ended-exploration-results.jpg}{Curvas de progreso de Voyager en las métricas de árbol tecnológico y de exploración abierta}{Diapositiva V037 recuperada de la grabación oficial; 00:30:21}

\begin{knowledgebox}{Lectura de la figura: primero el presupuesto del eje horizontal, después el significado de las curvas}
El eje horizontal representa el presupuesto de iteraciones o de exploración, y el eje vertical corresponde a objetos desbloqueados, hitos o task progress; primero se compara la velocidad de ascenso de Voyager y de las baselines con el mismo presupuesto, y después la meseta final. Las curvas respaldan que el ciclo memory/curriculum/code es más eficiente en Minecraft; no respaldan que «el modelo paramétrico ya adquirió inteligencia general» ni una extrapolación directa a robots reales.
\end{knowledgebox}

\lecturefigure{slide-38-map-coverage.jpg}{La comparación de cobertura del mapa muestra el alcance espacial de la exploración abierta}{Diapositiva V038 recuperada de la grabación oficial; 00:31:06}

Los colores y las regiones del mapa representan las trayectorias de exploración o el área cubierta por distintos agents. Primero se compara el área alcanzable y después se verifica que no se trate solo de una caminata aleatoria; solo en combinación con la finalización de tareas y la obtención de recursos se puede juzgar si la exploración es útil. Una mayor cobertura respalda que el agent visita activamente más estados, pero no indica que cada región se haya comprendido ni que haya aprendido un modelo causal de largo plazo.

\teachervoice{Al presentar los resultados, el ponente enfatiza que Voyager desbloquea la technology tree más rápido y explora un mapa más grande; al mismo tiempo, esto sigue siendo progress en Minecraft. Es importante situar las métricas dentro de los límites del entorno: la cobertura del mapa y la cantidad de objetos son proxies de exploración abierta, no una medición completa de la general intelligence.}

\subsection{Resumen de la sección}

Voyager usa code-as-action para comprimir comportamientos largos en programas ejecutables, los revisa repetidamente con el environment feedback y el critic, acumula executable memory con la skill library y decide la siguiente tarea con el automatic curriculum. Muestra una arquitectura de mejora continua que no requiere actualizar los pesos del modelo, pero su capacidad sigue limitada por la API, el critic, la memory hygiene, la task proposal y las métricas de Minecraft.
\section{Eureka: el LLM diseña el reward y el RL aprende el control de bajo nivel}

Voyager es adecuado para una symbolic action interface, pero el dexterous control requiere acciones continuas de alta frecuencia, y no es realista que el LLM escriba directamente cada motor command. Eureka cambia el nivel de la interfaz: el LLM lee el simulator environment code y genera un reward program; después, el RL interno aprende la policy mediante gradientes. La \term{reward reflection (reflexión sobre la recompensa)} consiste en que el modelo corrige el reward code a partir de las curvas de entrenamiento, el reward desglosado por términos y las conductas fallidas, en lugar de hacer una autoevaluación genérica de una respuesta en lenguaje natural.

\subsection{Hybrid-gradient Architecture}

Eureka divide el sistema en una code search externa de caja negra y una policy optimization interna de caja blanca. El LLM externo no necesita retropropagación; su función es proponer un reward ejecutable. El RL interno usa el muestreo en paralelo del simulator y el gradient update para convertir ese objective en acciones. Ambos niveles se conectan mediante las estadísticas de entrenamiento, de modo que el modelo de lenguaje puede «guiar» el control neuronal de bajo nivel sin producir directamente actions de alta frecuencia.

\lecturefigure{slide-39-eureka-architecture.jpg}{Eureka: arquitectura de dos niveles con el código del entorno, el reward program y la RL policy}{Diapositiva V039 recuperada de la grabación oficial; 00:33:36}

\begin{knowledgebox}{Lectura de la figura: primero hay que distinguir quién optimiza qué}
El LLM externo busca el reward code y el RL interno optimiza la policy con un reward fijo. El código del entorno expone al LLM la observation, el state y las magnitudes calculables; los resultados del entrenamiento regresan como un performance summary. La contribución de Eureka no es que GPT-4 aprenda por sí mismo el motor control, sino que reescribe la objective interface de un sistema de aprendizaje diferenciable.
\end{knowledgebox}

El proceso de dos niveles puede escribirse así:
\[
\phi_{k+1}\sim p_{\mathrm{LLM}}(\phi\mid \text{env code},\,\mathcal{H}_k),
\qquad
\theta_k^*=\arg\max_\theta J(\theta; r_{\phi_k}),
\qquad
\mathcal{H}_{k+1}=\mathcal{H}_k\cup\{\phi_k,\operatorname{metrics}(\theta_k^*)\}.
\]
Aquí $\phi_k$ es el reward program de la iteración $k$, $r_{\phi_k}$ es el reward que este calcula, $\theta_k^*$ son los parámetros de la policy que aprende el RL interno, $J$ es el expected return y $\mathcal{H}_k$ es el historial de rewards y retroalimentación del entrenamiento. El nivel externo modifica el código y el interno modifica los pesos; sus mecanismos de aprendizaje son distintos.

\teachervoice{El ponente llama a Voyager una no-gradient architecture y a Eureka una hybrid-gradient architecture: el ciclo externo de GPT-4 se encarga de elegir la reward function y el ciclo interno de la red neuronal aprende el control mediante el RL gradient. Esta distinción evita describir erróneamente «un sistema en el que participa un LLM» como «un sistema en el que todos los componentes se entrenan de extremo a extremo con el LLM».}

\subsection{Reward Reflection e In-context Evolution}

Que un reward code se ejecute no significa que induzca la conducta correcta. Por ello, Eureka condensa en retroalimentación cada reward term, la tasa de éxito del entrenamiento y la policy behavior, para que el LLM determine qué términos pesan demasiado, cuáles pesan muy poco o cuáles apuntan en la dirección equivocada, y luego genere código e hyperparameters nuevos. El objeto de la reflexión es una specification ejecutable y sus consecuencias.

\lecturefigure{slide-40-eureka-reward-reflection.jpg}{Eureka reflexiona sobre el diseño del reward a partir de la retroalimentación del entrenamiento}{Diapositiva V040 recuperada de la grabación oficial; 00:35:24}

\begin{knowledgebox}{Lectura de la figura: la retroalimentación debe permitir localizar el reward term}
La diapositiva contiene el reward code, las estadísticas de entrenamiento y el texto de reflexión. Primero se revisan la escala y el signo de cada term; luego, si la policy obtiene puntuaciones altas mediante posturas anómalas; por último, si la nueva sugerencia apunta a una observable failure. Si solo se dispone del retorno total, sin term-wise metrics, al LLM le resulta difícil distinguir entre «no se alcanzó el objetivo» y «un término de shaping dominó la optimización».
\end{knowledgebox}

\lecturefigure{slide-41-eureka-reward-code-diff.jpg}{El reward code y sus pesos evolucionan a lo largo de varias rondas de retroalimentación}{Diapositiva V041 recuperada de la grabación oficial; 00:35:36}

El diff del código muestra cómo la specification se modifica paso a paso. Primero se comparan los reward terms agregados o eliminados; después, los coefficients y la normalization. Lo importante no es que el código se alargue, sino que la conducta aprendida se acerque más a la intención de la tarea. Las mejoras de la figura respaldan la eficacia de la in-context reward evolution, pero no garantizan la misma robustez ante un simulator nuevo, un robot real o un adversarial state.

\begin{lstlisting}[caption={Pseudocódigo conceptual del ciclo externo de reward generation y el ciclo interno de entrenamiento con RL en Eureka}]
history = []
for generation in range(max_generations):
    candidates = llm.generate_reward_programs(env_code, history)
    evaluated = []
    for reward_program in candidates:
        policy, metrics = train_rl(env, reward_program)
        rollout_summary = evaluate_policy(policy, env)
        evaluated.append((reward_program, metrics, rollout_summary))
    best = select_valid_candidate(evaluated)
    reflection = llm.reflect(best, term_wise_metrics=True)
    history.append((best, reflection))
\end{lstlisting}

\begin{warningbox}{El LLM-generated reward code debe aislarse y verificarse}
Antes de ejecutarlo deben aplicarse un syntax/type check, una API allowlist, una verificación de rangos numéricos y un timeout; después del entrenamiento hay que revisar NaN, reward saturation, unsafe contact, constraint violation y held-out scenarios. La reward optimization busca activamente cualquier loophole, por lo que «el código se ejecuta» es solo el requisito mínimo.
\end{warningbox}

\subsection{Dexterity y Simulator Scaling}

Los resultados de Eureka se centran en la dexterous manipulation: tareas con posturas finas, contacto y control continuo, que difícilmente puede generar el LLM paso a paso de forma directa. El entrenamiento en paralelo a gran escala en el simulador Isaac permite que cada reward candidate reciba una feedback relativamente rápida y comparable; por lo tanto, el beneficio del algoritmo depende de contar con un environment code confiable y con recursos de cómputo suficientes.

\lecturefigure{slide-42-dexterity-scaling.jpg}{Las dexterous tasks y la escala del Isaac simulator sustentan la reward search}{Diapositiva V042 recuperada de la grabación oficial; 00:35:54}

\begin{knowledgebox}{Lectura de la figura: los resultados no pueden explicarse por separado de la infraestructura}
Las curvas de la izquierda comparan el desempeño en las tareas del reward learned/evolved con el human baseline; la parte derecha muestra el simulator scaling. Primero hay que confirmar qué indica el eje horizontal (tarea o escala de recursos) y el vertical (success/score); después, comparar si la ventaja se mantiene de forma estable entre tareas. La conclusión es que la LLM reward search puede producir un objective sólido dentro de un simulator dado; no demuestra que estén resueltos la simulation-to-real gap, el ruido de los sensores ni la seguridad del hardware.
\end{knowledgebox}

\teachervoice{El ponente destaca que Eureka supera al human-designed reward en varias dexterous tasks, y señala a la vez que la escala del simulator hace viable la búsqueda externa. Dicho de otro modo, la capacidad del método surge de la acción conjunta del LLM prior, el environment source code, el RL optimizer y la infraestructura en paralelo; no puede atribuirse solo al prompt.}

\subsection{Resumen de la sección}

Eureka coloca al LLM en el nivel de objective-design y deja el RL basado en gradientes en el nivel de control continuo. La reward reflection modifica una specification ejecutable a partir de la retroalimentación real del entrenamiento, con lo que forma una hybrid architecture de code search externa y policy learning interno. Establece una interfaz entre el reasoning de alto nivel y el motor control de bajo nivel, y a la vez exige un sandbox estricto, term-wise metrics, constraint audit y límites claros del simulator.
\section{Qué aprender del video a escala de Internet: representation, reward y action}

Después de MineCLIP, Voyager y Eureka, la clase pasa a direcciones de investigación de más largo plazo. El video de Internet es barato, dinámico y abarca una gran cantidad de conducta humana, pero proviene de perspectivas en tercera persona o no controladas y por lo general carece de etiquetas de agent action, state o grounding. El problema no debe llamarse de forma genérica «imitar a partir de video»; conviene descomponerlo en tres objetivos de supervisión: aprender una representation, aprender una reward, o inferir la action para luego hacer behavior cloning.

\subsection{Perspectiva: dos rutas hacia el foundation agent}

El ponente divide las direcciones futuras en dos tipos: seguir aprendiendo de internet-scale video y construir un multimodal foundation model nativo. La primera amplía los priors de conducta; la segunda unifica las interfaces de texto, imagen, video, sonido y action. Ambas deben volver, al final, al active loop y someterse a la retroalimentación de la embodiment y del environment.

\lecturefigure{slide-43-looking-forward-title.jpg}{Transición de los sistemas actuales a la dirección de largo plazo del generalist agent}{Diapositiva V043 recuperada de la grabación oficial; 00:37:36}

Esta página funciona como portada de sección: los sistemas anteriores ya demostraron que reward, code y memory son viables; lo que sigue es preguntar cómo ampliar los datos y las modalidades. Al leer la figura no hay que extrapolar directamente el éxito de chat/image/music a los agents: los datos de acción contienen causalidad, control y riesgo, y no se obtienen automáticamente con más tokens estáticos.

\lecturefigure{slide-44-two-foundation-directions.jpg}{Internet-scale video learning y multimodal foundation models}{Diapositiva V044 recuperada de la grabación oficial; 00:38:09}

\begin{importantbox}{Las dos rutas convergen en el action grounding}
El video learning aporta conducta dinámica y estrategias humanas; el multimodal model aporta una token/interface unificada. Aun así, el sistema necesita saber «qué action ejecuta quién y qué consecuencias tiene para la embodiment actual». Sin active grounding, el modelo aprende, como mucho, una observational correlation.
\end{importantbox}

\subsection{Ventajas y carencias del video de Internet}

El video está más cerca de la conducta que el texto, porque muestra cómo cambian los objetos con el tiempo, cómo la mano entra en contacto con las herramientas y cómo se desarrolla una tarea. Sin embargo, quien observa no sabe qué ocurre fuera de cuadro, ni ve el action command preciso, la fuerza, el tacto o los reintentos tras un fallo. La \term{behavior cloning (clonación de comportamiento)} entrena una policy con supervisión a partir de los pares observation--action de las demostraciones; cuando no hay action label, primero hay que recuperar o aproximar esos pares.

\lecturefigure{slide-45-internet-video-properties.jpg}{Internet video: barato, dinámico y rico, pero sin action ni grounding}{Diapositiva V045 recuperada de la grabación oficial; 00:38:48}

\begin{knowledgebox}{Términos clave: tres rutas para aprender del video}
\begin{tabularx}{\textwidth}{@{}L{0.19\textwidth}X X@{}}
\toprule
Ruta & Qué aprende & Carencia principal \\
\midrule
Representation & Representaciones espaciotemporales y semánticas transferibles & Una buena representación no implica saber controlar \\
Reward learning & Similitud/valor entre goal y trajectory & Reward bias y shortcut \\
Action / behavior cloning & Predecir una action ejecutable a partir de la observation & El video suele carecer de etiquetas de acción y la embodiment es distinta \\
\bottomrule
\end{tabularx}
\end{knowledgebox}

\lecturefigure{slide-46-internet-video-example.jpg}{El video de Internet en tercera persona muestra la conducta, pero oculta las señales de control}{Diapositiva V046 recuperada de la grabación oficial; 00:39:39}

En este ejemplo conviene mirar primero la secuencia temporal entre la mano, el objeto y el resultado, y después preguntar dónde está la action real. Los píxeles solo muestran el movimiento superficial; no proporcionan comandos articulares, fuerza aplicada, intención ni el estado fuera de cuadro, y una misma trayectoria visual puede ser producida por distintos controles. El video sirve para aprender semantic dynamics, pero requiere supuestos adicionales para convertirse en motor supervision de un robot concreto.

\teachervoice{El ponente describe el video de Internet como «cheap, dynamic, rich in human behavior» e inmediatamente añade que un observador pasivo no obtiene action labels ni grounding. Este giro es importante: la escala de los datos resuelve la coverage, pero no resuelve automáticamente el causal control; aquí reaparece el problema del gatito activo.}

\subsection{Ruta uno: Video Representation Learning}

La primera ruta aprende primero una representation espaciotemporal y después aplica fine-tune con una pequeña cantidad de agent data específica del dominio. Mediante future prediction, masked modeling o video--text alignment, el modelo puede obtener representaciones de objetos, fases de acción y cambios de escena, pero la capacidad de control sigue dependiendo de la action data y de la interface posteriores.

\lecturefigure{slide-47-representation-from-video.jpg}{Preentrenar la representation con video y luego adaptarla a tareas del agent}{Diapositiva V047 recuperada de la grabación oficial; 00:40:18}

\begin{knowledgebox}{Lectura de la figura: el preentrenamiento y el ajuste para control cumplen funciones distintas}
La etapa de video a gran escala de la figura aporta rasgos dinámicos generales; solo la pequeña cantidad posterior de agent data conecta la representation con un action space concreto. Primero compara el encoder congelado, el fine-tune y el entrenamiento desde cero, y luego revisa el tamaño de los datos posteriores; si no hay una action-grounded evaluation, solo se puede afirmar que la feature es útil, no que la policy haya aprendido la tarea.
\end{knowledgebox}

\subsection{Ruta dos: aprender la reward a partir del video}

La ruta anterior solo exige que el video preentrene una representación útil; esta sección va más allá y hace que el video participe directamente en la optimización de la tarea: MineCLIP convierte «cuánto se parece la trayectoria al objetivo» en una reward que guía al agent a explorar acciones en su propio entorno. No exige que el video y el robot compartan un control preciso, por lo que cruza interfaces con más facilidad que la imitación acción por acción; el costo es que hay que verificar si la semántica visual se transfiere y evitar que la policy explote activamente los puntos ciegos del reward model.

\lecturefigure{slide-48-reward-from-video.jpg}{Las representaciones video-language y goal-conditioned pueden formar una learned reward}{Diapositiva V048 recuperada de la grabación oficial; 00:40:45}

\begin{knowledgebox}{Lectura de la figura: la reward transfer es más débil, pero más flexible, que la action transfer}
El video solo necesita indicarle al sistema «si el resultado se ve más cerca del objetivo»; la action concreta la explora la policy en su propia embodiment, por lo que puede cruzar interfaces de control. Pero cuando la reward solo restringe la outcome appearance, puede pasar por alto la trayectoria, la seguridad o condiciones implícitas. Al leer los resultados hay que revisar a la vez el final state, la trajectory y la constraint.
\end{knowledgebox}

\subsection{Ruta tres: recuperar la action y hacer behavior cloning}

La tercera ruta es la más cercana al imitation learning: inferir una latent action a partir del video o usar un modelo para etiquetar la action, y luego entrenar una behavior policy. Video PreTraining (VPT), en Minecraft, entrena un inverse dynamics model con una pequeña cantidad de datos con etiquetas de teclado y ratón, y después completa la action de una gran cantidad de video sin etiquetar; este método depende de que la perspectiva de observación, el espacio de acciones y la embodiment coincidan lo suficiente.

\lecturefigure{slide-49-behavior-cloning-from-video.jpg}{Recuperar action labels a partir del video y luego hacer behavior cloning}{Diapositiva V049 recuperada de la grabación oficial; 00:42:33}

Si $z_t$ denota la latent action inferida a partir de cuadros adyacentes y la agent action real es $a_t$, puede escribirse:
\[
z_t\sim q_\psi(z\mid o_t,o_{t+1}),\qquad
a_t\sim p_\omega(a\mid z_t,e),\qquad
\mathcal{L}_{\mathrm{BC}}=-\sum_t\log\pi_\theta(a_t\mid o_{\le t},g).
\]
Aquí $q_\psi$ es el inverse/latent-action model, $p_\omega$ mapea la latent action a una acción ejecutable de la embodiment $e$, $\pi_\theta$ es la behavior policy y $g$ es el objetivo. Si la embodiment, la camera o el action support del video y del agent no coinciden, $p_\omega$ se convierte en la principal fuente de error.

\begin{warningbox}{Observation similarity no equivale a action identifiability}
A partir de solo dos imágenes por lo general no es posible recuperar de forma única la fuerza aplicada, la velocidad ni los contactos ocultos; el video en tercera persona además presenta movimiento de cámara y oclusiones. La action reconstruction debe reportar su uncertainty y permitir que la downstream policy corrija mediante la interacción con el entorno, en lugar de tratar las pseudoetiquetas como ground truth.
\end{warningbox}

\teachervoice{El ponente separa explícitamente representation, reward y action learning en tres rutas para evitar que «entrenar agents con video» se vuelva un lema vacío. La necesidad de action labels crece de una a otra: representation es la más débil, reward queda en medio y behavior cloning es la más fuerte; el sistema debe elegir la interfaz según la supervisión que pueda obtener.}

\subsection{Resumen de la sección}

El internet video ofrece datos de conducta humana baratos, dinámicos y de amplia cobertura, pero carece de action labels y grounding. El representation learning aprende rasgos, el reward learning aprende la similitud con el objetivo y la behavior cloning intenta recuperar las acciones; las tres rutas difieren en intensidad de supervisión, supuestos de transferencia y modos de falla. La escala puede ampliar los priors, pero la active environment feedback sigue siendo indispensable para verificar la causalidad de las acciones.
\section{Multimodal Foundation Agents: VIMA, RT-2 y RoboCat}

La ruta del video amplía los datos; la otra ruta unifica la interfaz. \term{multimodal prompting (prompting multimodal)} se refiere a que en el prompt se intercalan texto, imágenes, fotogramas objetivo, demostraciones en video u otros tokens perceptuales, y el modelo produce una action a partir de ellos. No consiste en añadir una imagen decorativa a un prompt de lenguaje, sino en hacer que los objetos visuales, las restricciones y las demostraciones formen parte directa de la especificación de la tarea.

\subsection{Unificar Text, Image, Video, Audio y Action}

Un multimodal foundation model intenta proyectar entradas distintas a un espacio de secuencias compartido y luego generar lenguaje o action según el contexto. Para un agent, la ventaja es que la tarea puede expresarse como «pon este objeto en la posición que muestra esta imagen» o mediante «una demostración»; el reto es que cada modality tiene requisitos distintos de tokenización, temporalidad, precisión y safety.

\lecturefigure{slide-50-multimodal-foundation-model.jpg}{Visión de un foundation model que unifica texto, imagen, video, sonido y action}{Diapositiva V050 recuperada de la grabación oficial; 00:43:03}

\begin{knowledgebox}{Lectura de la figura: un token space unificado no equivale a un significado físico unificado}
En la figura, todas las modality entran en un solo modelo, pero el costo de un error es distinto para un token de texto, un patch visual y una robot action. Primero observe cómo se codifica la entrada y después si la salida está sujeta a una action-space constraint; una architecture unificada favorece la knowledge transfer, pero sigue necesitando encoder específicos de cada modality, controller y safety validation.
\end{knowledgebox}

\lecturefigure{slide-51-multimodal-prompting-robots.jpg}{Una tarea robótica puede definirse con un prompt que intercala elementos visuales y de lenguaje}{Diapositiva V051 recuperada de la grabación oficial; 00:43:48}

El prompt puede incluir imágenes del objeto objetivo, esquemas de posición, zonas prohibidas y relaciones en lenguaje natural. Al leer la figura, primero distinga «a qué objeto se refiere», «cuál es el estado objetivo» y «qué constraint debe cumplir la acción», y después observe la robot trajectory. El prompt multimodal aumenta la specification bandwidth, pero también introduce ambiguity; el modelo debe hacer grounding en la scene actual, en lugar de emparejar plantillas de apariencia vistas en el entrenamiento.

\subsection{VIMA: Visual Manipulation with Multimodal Prompts}

VIMA introduce en un Transformer tokens de objetos visuales intercalados con tokens de texto y predice la robot action de forma autoregressive. Su aporte principal no es un manipulation benchmark aislado, sino expresar con un mismo model/interface familias de tareas distintas, como visual goal reaching, novel concept grounding, one-shot demonstration y constraint satisfaction.

\lecturefigure{slide-52-vima-architecture.jpg}{VIMA modela de forma unificada el multimodal prompt, la observation y la action}{Diapositiva V052 recuperada de la grabación oficial; 00:44:33}

\begin{knowledgebox}{Lectura de la figura: la conexión entre el prompt encoder y el action decoder}
Primero observe cómo las imágenes y el texto del prompt forman la token sequence; luego, cómo la observation actual se incorpora al contexto, y por último, cómo el action decoder produce tokens de pose o de control. El Transformer compartido de VIMA favorece la transferencia entre tareas, pero el tabletop world, el object set y la action representation del benchmark siguen limitando el alcance de su evidencia.
\end{knowledgebox}

Si los prompt tokens son $p_{1:m}$, las observation históricas son $o_{1:t}$ y la action sequence es $a_{1:t}$, el autoregressive objective puede escribirse como:
\[
\mathcal{L}_{\mathrm{VIMA}}(\theta)
=-\sum_{t=1}^{T}\log p_\theta\!\left(a_t\mid p_{1:m},o_{1:t},a_{<t}\right).
\]
Aquí $p_{1:m}$ puede intercalar texto y objetos visuales, $o_{1:t}$ son las observation actual e históricas, $a_{<t}$ son las acciones previas, $T$ es el episode horizon y $\theta$ son los parámetros del modelo. La fórmula expresa un conditional sequence modeling unificado, pero no implica que todo action error equivalga a un error en un token de lenguaje.

\subsection{Límites de la evidencia para las cuatro capacidades de Prompt}

Las cuatro diapositivas siguientes evalúan, respectivamente, objetivos visuales, conceptos temporales, una sola demostración y restricciones. Deben verse como formas distintas de task specification, no como cuatro demo sin relación entre sí; al leer las figuras, verifique continuamente qué información aporta el prompt, sobre qué objeto debe hacer grounding el modelo y si la evaluación se realiza solo dentro de una tabletop distribution controlada.

\lecturefigure{slide-53-vima-visual-goal-reaching.jpg}{VIMA: una imagen objetivo especifica el estado visual que se desea alcanzar}{Diapositiva V053 recuperada de la grabación oficial; 00:44:51}

El objetivo visual reduce la carga de describir relaciones espaciales con lenguaje. Primero compare la scene inicial con la image objetivo y luego verifique las posiciones relativas de los objetos después de la acción; el éxito indica que el modelo puede convertir la apariencia objetivo en una manipulation sequence, pero por sí solo no demuestra que comprenda la causalidad física en general ni las oclusiones del mundo real.

\lecturefigure{slide-54-vima-concept-grounding.jpg}{VIMA: grounding de conceptos nuevos definidos temporalmente dentro del prompt}{Diapositiva V054 recuperada de la grabación oficial; 00:45:03}

Palabras sin significado establecido, como «wug» y «blicket», se definen temporalmente mediante su correspondencia visual en el prompt. Al leer la figura, primero identifique el binding entre token y objeto, y luego vea si el modelo traslada la relación a la action; esto respalda el in-context concept grounding, no demuestra que el modelo haya visto el significado de esas palabras durante el preentrenamiento.

\lecturefigure{slide-55-vima-one-shot-demonstration.jpg}{VIMA: inferir la tarea a partir de una sola demostración visual y reproducirla}{Diapositiva V055 recuperada de la grabación oficial; 00:45:15}

Una sola demostración aporta una trajectory-level specification. Primero observe los cambios de los objetos en la demonstration y luego si la test scene tiene posiciones o instancias nuevas; el éxito exige que el modelo extraiga relaciones en lugar de copiar coordenadas absolutas. Que una sola demo pueda generalizarse sigue dependiendo de la similitud de la familia de tareas, la cámara y la object distribution.

\lecturefigure{slide-56-vima-constraint-satisfaction.jpg}{VIMA: completar una manipulation bajo restricciones visuales}{Diapositiva V056 recuperada de la grabación oficial; 00:45:48}

El constraint prompt le indica al agent qué zona, objeto o trayectoria no debe tocar. Al leer la figura, no basta con verificar el objetivo final: también hay que comprobar si la trajectory viola restricciones intermedias. Esto muestra que la agent evaluation debe incluir la seguridad del proceso, y no solo el final frame. Las restricciones en una mesa controlada siguen siendo distintas de la force, la collision y la human safety de un robot real.

\teachervoice{El ponente dice que VIMA «tokenizes everything» y que, con interleaved image--text tokens, unifica tareas que antes requerían un pipeline especializado. Esto debe entenderse como una unificación de la modeling interface, no como la desaparición del problema de control físico; el grounding visual, el action decoder y las restricciones del entorno siguen asumiendo cada uno su propia responsabilidad.}

\subsection{RT-2 y RoboCat: ampliar los Robot Data y la Transfer}

Al final, el curso usa RT-2 y RoboCat para mostrar ejemplos posteriores de la misma recipe. RT-2 hace co-fine-tune del web-scale vision-language knowledge junto con robot demonstration, y representa la action como tokens que pueden generarse; RoboCat entrena una policy unificada que abarca distintas tareas y robot embodiment. Ambos amplían la evidencia de transfer, pero siguen dependiendo de la mezcla de datos concreta y del action support.

\lecturefigure{slide-57-rt2-robocat.jpg}{RT-2 y RoboCat extienden la ruta del multimodal generalist robot}{Diapositiva V057 recuperada de la grabación oficial; 00:46:12}

\begin{knowledgebox}{Términos clave: tres interfaces de generalist-agent}
\begin{tabularx}{\textwidth}{@{}L{0.17\textwidth}L{0.25\textwidth}X@{}}
\toprule
Sistema & Interfaz principal & Mecanismo central y límites \\
\midrule
Voyager & JavaScript code-as-action & Horizon largo, componible; depende de una symbolic API \\
Eureka & LLM-generated reward code & Aprende control continuo con RL; depende del simulator y de la reward audit \\
VIMA / RT-2 / RoboCat & De multimodal token a robot action & Representación compartida y transferencia de datos; depende del action grounding y de la embodiment coverage \\
\bottomrule
\end{tabularx}
\end{knowledgebox}

\teachervoice{El ponente explica RT-2 como un modelo preentrenado primero con internet-scale data y luego sometido a fine-tune con human-collected robot demonstrations; RoboCat, por su parte, muestra una sola policy que abarca distintos robots y tareas. En clase se presentan como un promising follow-up, sin afirmar que la generalist robotics ya esté resuelta.}

\subsection{Resumen de la sección}

El multimodal prompting amplía la specification de la tarea: texto, imágenes objetivo, conceptos temporales, demostraciones y constraint pueden entrar en un contexto unificado. VIMA demuestra que varias clases de tabletop task pueden expresarse con un mismo Transformer/action interface, y RT-2 y RoboCat amplían además el web prior, los robot data y la embodiment transfer. Un modelo unificado aporta representaciones compartidas, pero el grounding, el control, la safety y el evidence boundary todavía deben verificarse por separado.
\section{Síntesis y ampliación: un Generalist Agent es un sistema de ciclo cerrado, no un solo modelo}

Toda la clase parte del gatito activo y termina en una descomposición del sistema que puede trasladarse a otros casos. El entorno abierto genera tareas que no se agotan y el conocimiento de internet aporta conocimiento previo humano. El foundation model procesa el contexto de lenguaje y multimodal. La reward/code/action interface conecta ese conocimiento previo con la conducta, el critic y el environment feedback evalúan las consecuencias, y la skill memory y el curriculum sostienen la acumulación a largo plazo. Cualquier benchmark puntual solo puede validar una parte de este sistema.

\subsection{Cierre del curso: MineDojo, Agent Learning y Multimodal Prompts}

Las tres últimas diapositivas repasan los artículos, las demos de agentes y los multimodal prompts, y sitúan estos trabajos dentro de un objetivo común. No forman una «lista de productos» en la que uno sustituye a otro. MineDojo aporta el environment/data substrate, MineCLIP aprende una reward semántica, Voyager aprende una memory ejecutable, Eureka busca reward programs, y VIMA/RT-2/RoboCat exploran un action model multimodal unificado.

\lecturefigure{slide-58-minedojo-paper-overview.jpg}{Repaso del artículo de MineDojo: infraestructura de entorno abierto y de conocimiento de internet}{Diapositiva V058 recuperada de la grabación oficial; 00:46:24}

\begin{knowledgebox}{Lectura de la figura: reducir el collage a las responsabilidades del sistema}
La diapositiva muestra a la vez capturas del entorno, material de internet e información del artículo. Primero identifique las tres capas: environment, knowledge database y benchmark. Después pregúntese qué evidence aporta cada una. La diapositiva demuestra que la plataforma cubre mucho, pero el collage no basta para juzgar la policy quality. La validación real viene del protocolo de tareas, la reward, el baseline y artifacts reproducibles.
\end{knowledgebox}

\lecturefigure{slide-59-agent-learning-examples.jpg}{Ejemplos de conductas y tareas de agent learning en mundo abierto}{Diapositiva V059 recuperada de la grabación oficial; 00:46:27}

Los ejemplos muestran que el agente puede recolectar, fabricar, explorar o interactuar en escenarios distintos. Al leer la figura, no sume varias capturas exitosas como si fueran una success rate. Las capturas ilustran la task diversity. La capacidad cuantitativa debe medirse con held-out evaluation, failure distribution y presupuesto. En un mundo abierto es especialmente necesario informar lo que el agente no logró hacer.

\lecturefigure{slide-60-multimodal-prompts-summary.jpg}{El multimodal prompt unifica metas visuales, relaciones lingüísticas y tareas robóticas}{Diapositiva V060 recuperada de la grabación oficial; 00:46:51}

Esta diapositiva vuelve a la interfaz central de VIMA: el prompt no es un instruction string fijo, sino un task program multimodal que se puede componer. Primero observe los tipos de token y el binding de objetos. Después observe la action que produce el modelo. Esto permite specification flexibility, pero también aumenta la responsabilidad ante la prompt ambiguity, el adversarial input y la safety validation.

\subsection{Grand Challenge: Human-level Minecraft AI}

El curso cierra con Minecraft de nivel humano como reto para la comunidad. Aquí «human-level» no puede reducirse a vencer a un jefe o reunir un conjunto de objetos. Debe incluir comprensión del lenguaje, memoria a largo plazo, construcción creativa, colaboración social, recuperación tras fallas, adaptación entre versiones y aprendizaje continuo con metas nuevas. Es una línea de investigación, no un logro que la clase afirme haber alcanzado.

\lecturefigure{slide-61-human-level-minecraft-challenge.jpg}{Grand Challenge: Human-level Minecraft AI}{Diapositiva V061 recuperada de la grabación oficial; 00:47:18}

\begin{warningbox}{No deje que un leaderboard sustituya a «Human-level»}
La velocidad en una prueba, la cobertura del mapa o la cantidad de objetos solo miden capacidades parciales. Una evaluación más completa informa por separado task proposal, knowledge use, planning horizon, memory retention, creative outcome, robustness, safety y human collaboration, y hace pruebas con seeds, versiones y metas no vistos.
\end{warningbox}

\teachervoice{Al final, el ponente deja el Minecraft de nivel humano como tarea para toda la comunidad y comenta medio en broma que también es un pretexto para «jugar Minecraft en el trabajo». La enseñanza de fondo es que el trabajo actual solo toca la superficie de los agentes en mundo abierto. La siguiente etapa necesita mejor aprendizaje a partir de video, modelos fundacionales multimodales más potentes y ciclos cerrados activos. Un conjunto de demos vistosas no es el punto de llegada.}

\subsection{Lista de verificación del sistema en ocho dimensiones}

Para trasladar esta clase a otros proyectos de agentes, cualquier sistema puede auditarse punto por punto en las ocho dimensiones siguientes. Cubren problem, data, model, interface, learning loop y evidence. Así se evita que un equipo confunda la fluidez lingüística del foundation model con una capacidad completa de acción.

\begin{longtable}{@{}L{0.18\textwidth}L{0.34\textwidth}L{0.36\textwidth}@{}}
\toprule
Dimensión & Pregunta central & Falla típica \\
\midrule
Environment & ¿El mundo tiene reglas que se pueden combinar y estado a largo plazo? & Se llama mundo abierto a un benchmark cerrado \\
Goal & ¿Cómo se plantea la meta, cómo se elimina su ambigüedad y cómo se decide que se cumplió? & Un reward shortcut aprovecha un prompt ambiguo \\
Knowledge & ¿Qué conocimiento previo viene de la web, de video, de personas o del simulador? & Fuentes con ruido que no corresponden al estado actual \\
Interface & ¿La salida es action, code, reward o skill? & El nivel de abstracción no corresponde a la dinámica de la tarea \\
Feedback & ¿Quién aporta success, error y constraint que puedan verificarse? & El critic se valida a sí mismo; reward sin calibrar \\
Memory & ¿Cómo se almacena, recupera, actualiza y da de baja la experiencia exitosa? & Una skill errónea contamina por mucho tiempo las tareas siguientes \\
Curriculum & ¿Cómo equilibra la siguiente tarea factibilidad, novedad y riesgo? & Demasiado difícil y se estanca; demasiado fácil y se repite; solo persigue la reward \\
Evidence & ¿Qué demuestra la métrica y qué no demuestra? & Una ventaja local en el simulador se extrapola a inteligencia general \\
\bottomrule
\end{longtable}

\begin{importantbox}{La clase en una sola frase}
Un Generalist agent no es «un modelo más grande». Conecta en un ciclo cerrado \textbf{el entorno abierto, el conocimiento previo de internet, las interfaces ejecutables, la retroalimentación real, la memoria persistente y el currículo automático}. El valor del modelo depende de que, dentro de ese ciclo, convierta passive knowledge en active competence verificada.
\end{importantbox}

\subsection{Los cuatro criterios de ingeniería que más vale conservar}

Aunque esta clase usa Minecraft y robots como ejemplos, los cuatro criterios se trasladan directamente a tool-using agents, browser agents, coding agents y scientific agents. Los cuatro sostienen que la capacidad de un sistema la determinan environment/interface/evidence, no el nombre del modelo por sí solo.

\begin{enumerate}
  \item \textbf{Elija primero la action abstraction.} El motor control de alta frecuencia, las API call, los programas y la reward specification requieren interfaces distintas. No asigne al modelo responsabilidades que no correspondan a su escala de tiempo.
  \item \textbf{Tome el execution result como truth source.} La autoevaluación en lenguaje solo sirve de apoyo. El success real debe venir del estado del entorno, las pruebas, los constraint y la revisión humana.
  \item \textbf{Guarde la memory como artifacts verificables.} El código, las skill, las trajectory y los reward program deben incluir precondiciones, versión, pruebas y una política de caducidad.
  \item \textbf{Incluya en las conclusiones los límites de las métricas.} Las curvas y las demos solo respaldan afirmaciones dentro de su protocolo de evaluación. La capacidad entre entornos, entre embodiment y en seguridad requiere evidencia independiente.
\end{enumerate}

\subsection{Lecturas adicionales}

\begingroup
\footnotesize
\begin{itemize}[itemsep=0pt,topsep=0.15em,parsep=0pt]
  \item Fan et al., \href{https://arxiv.org/abs/2206.08853}{\emph{MineDojo: Building Open-Ended Embodied Agents with Internet-Scale Knowledge}}
  \item Wang et al., \href{https://arxiv.org/abs/2305.16291}{\emph{Voyager: An Open-Ended Embodied Agent with Large Language Models}}
  \item Ma et al., \href{https://arxiv.org/abs/2310.12931}{\emph{Eureka: Human-Level Reward Design via Coding Large Language Models}}
  \item Jiang et al., \href{https://arxiv.org/abs/2210.03094}{\emph{VIMA: General Robot Manipulation with Multimodal Prompts}}
  \item Baker et al., \href{https://arxiv.org/abs/2206.11795}{\emph{Video PreTraining (VPT): Learning to Act by Watching Unlabeled Online Videos}}
  \item Brohan et al., \href{https://arxiv.org/abs/2307.15818}{\emph{RT-2: Vision-Language-Action Models Transfer Web Knowledge to Robotic Control}}
  \item Reed et al., \href{https://arxiv.org/abs/2306.11706}{\emph{RoboCat: A Self-Improving Foundation Agent for Robotic Manipulation}}
\end{itemize}

Al leer estos trabajos, se puede aplicar una y otra vez la lista de verificación de ocho dimensiones. Además, conviene preguntar si los datos de entrenamiento incluyen action consequence y si la evaluation abarca los failure. También conviene preguntar si la memory acumulará errores, si la reward podría aprovecharse con atajos y si la interfaz conviene al embodiment actual. Por último, hay que revisar si las conclusiones rebasan los límites del simulador, del conjunto de tareas y del horizonte de tiempo.
\endgroup

\end{document}
